\section{图表与评价}
\begin{frame}{Scaling curve：串行比例限制资源扩展收益}
\begin{columns}[T,onlytextwidth]
\begin{column}{.68\textwidth}\fig[height=4.55cm,keepaspectratio]{amdahl}\end{column}
\begin{column}{.28\textwidth}
\subhead{先读趋势，再读边界}
更高的并行比例带来更高的极限加速比。
\par\vspace{3mm}\metric{\linewidth}{$20\times$}{$f=0.95$ 时，$p\to\infty$ 的理论上限}
\par\vspace{3mm}\notebox{曲线由 Amdahl 模型计算，假设额外开销为零。}
\end{column}
\end{columns}
\insight{图题表达结论；坐标和图注定义数值含义；正文解释对研究问题的影响。}
\end{frame}

\begin{frame}{Roofline：让两个约束出现在同一张图中}
\begin{columns}[T,onlytextwidth]
\begin{column}{.65\textwidth}\fig{roofline}\end{column}
\begin{column}{.31\textwidth}
\subhead{归一化的解析关系}
\[r=\min(1,x)\]
\[x=\frac{BI}{P_{\max}},\qquad r=\frac{P}{P_{\max}}\]
\smalltext $B$：带宽；$I$：算术强度；$P_{\max}$：计算上限。
\par\vspace{3mm}\notebox{图中仅展示模型约束。具体设备的上限需使用相应精度、指令类型和测量口径。}
\end{column}
\end{columns}
\captiontext{所有轴均为无量纲归一化量；拐点 $x=1$ 对应带宽约束与计算约束相交。}
\end{frame}

\begin{frame}{数值表：同一模型中的收益与效率}
\renewcommand{\arraystretch}{1.10}
\begin{tabularx}{\textwidth}{@{}XXXX@{}}
\toprule 资源数量 $p$ & 归一化时间 $T_p/T_1$ & Speedup $S(p)$ & Efficiency $S(p)/p$ \\\midrule
1  & 1.0000 & 1.000 & 100.00\% \\
2  & 0.5250 & 1.905 & 95.24\% \\
4  & 0.2875 & 3.478 & 86.96\% \\
8  & 0.1688 & 5.926 & 74.07\% \\
16 & 0.1094 & 9.143 & 57.14\% \\
32 & 0.0797 & 12.549 & 39.22\% \\\bottomrule
\end{tabularx}
\vspace{2mm}\captiontext{解析计算：$f=0.95$，$T_p/T_1=0.05+0.95/p$。各列从未舍入数值独立取舍。}
\insight{资源增加时，Speedup 与资源利用效率呈现不同变化；评价需要同时说明二者。}
\end{frame}

\begin{frame}{实验设计：把比较条件写到展示内容中}
\begin{columns}[T,onlytextwidth]
\begin{column}{.48\textwidth}
\subhead{控制条件}
\begin{tabularx}{\linewidth}{@{}lX@{}}
\toprule 条件 & 记录内容 \\\midrule
输入 & 规模、分布、边界与随机种子 \\
软件 & Commit、Compiler、编译选项 \\
硬件 & GPU 型号、显存、运行状态 \\
计时 & Warm-up、同步位置、重复次数 \\\bottomrule
\end{tabularx}
\end{column}
\begin{column}{.48\textwidth}
\subhead{输出证据}
\begin{denseitems}
\item 正确性：参考结果与边界输入检查。
\item 性能：端到端延迟、吞吐与离散程度。
\item 机制：阶段计时、资源与访问指标。
\end{denseitems}
\notebox{性能测量与机制分析分别记录运行配置，检查分析工具对执行时间的影响。}
\end{column}
\end{columns}
\vspace{3mm}\captiontext{本页给出实验设计结构；在自己的报告中填入已执行的配置与结果文件位置。}
\end{frame}
